% SOURCE_AUTHORITY: sga5_fr_workpass.tex, lines 3961--4013 (LF-normalized unit).
% SOURCE_SHA256: 791F4EFFC5E02832D5D77ED03518C8156D6F07E4C8238B03545DB93D883FBB28
\subsection*{6.9}

Sejam $X$ (resp. $Y$) um $S$-esquema próprio e suave de dimensão relativa $m$ (resp. $n$), $Z=X\times_S Y$, e sejam $a:A\hookrightarrow Z$, $b:B\hookrightarrow Z$ imersões regulares tais que $a_2:A\to Y$ e $b_1:B\to X$ sejam finitos e planos (denotamos por $p_1:Z\to X$, $p_2:Z\to Y$ as projeções, e $a_i=p_i a$, $b_i=p_i b$), $c:C=A\times_Z B\to Z$, $L\in\ob D(X)_{\parf}$ e $M\in\ob D(Y)_{\parf}$ (o índice parf significa ``perfeito em sentido absoluto'', o que equivale também a perfeito relativamente a $S$, pois $X$ e $Y$ são suaves). Supõe-se que $C$ seja finito e plano sobre $S$. As hipóteses impostas implicam, portanto, que $c$ é uma imersão regular e que $a$ e $b$ são Tor-independentes.

Temos $p_1^*\Omega^1_{X/S}=\Omega^1_{Z/Y}$, $p_2^*\Omega^1_{Y/S}=\Omega^1_{Z/X}$, e uma decomposição
\[
\Omega^1_{Z/S}=\Omega^1_{Z/Y}\oplus\Omega^1_{Z/X},
\]
de onde, para todo inteiro $r$, obtém-se uma decomposição
\begin{equation}
\tag{6.9.0}
\Omega^r_{Z/S}=\bigoplus_{i+j=r}\Omega^i_{Z/Y}\otimes\Omega^j_{Z/X}.
\end{equation}
Por outro lado, de acordo com 6.1.2, temos $p_2^!M=p_2^*M\otimes\Omega^m_{Z/Y}[m]$, e, portanto,
\[
H^0(Z,a_*a^!\uRHom(p_1^*L,p_2^!M))=\Ext^m_{\OO_Z}(a_1^*L,a_2^*M\otimes\Omega^m_{Z/Y}),
\]
de onde, por 6.4.4 e 6.8.2, obtém-se um isomorfismo
\begin{equation}
\tag{6.9.1}
\Hom(a_1^*L,a_2^!M)=H^0(A,\uRHom(a_1^*L,a_2^*M)\otimes\Lambda^m N_{A/Z}^{\vee}\otimes\Omega^m_{Z/Y}).
\end{equation}
Da mesma forma, temos um isomorfismo
\begin{equation}
\tag{6.9.2}
\Hom(b_2^*M,b_1^!L)=H^0(B,\uRHom(b_2^*M,b_1^*L)\otimes\Lambda^n N_{B/Z}^{\vee}\otimes\Omega^n_{Z/X}).
\end{equation}
Sejam $u\in\Hom(a_1^*L,a_2^*M)$ e $v\in\Hom(b_2^*M,b_1^*L)$; denotemos por
\begin{equation}
\tag{6.9.3}
\langle u,v\rangle\in H^0(C,\OO_C)
\end{equation}
a seção $\operatorname{Tr}(u_Cv_C)=\operatorname{Tr}(v_Cu_C)$, onde $u_C\in\Hom(c_1^*L,c_2^*M)$ e $v_C\in\Hom(c_2^*M,c_1^*L)$ são induzidos por $u$ e $v$. Consideremos a correspondência cohomológica $u^!\in\Hom(a_1^*L,a_2^!M)$ (resp. $v^!\in\Hom(b_2^*M,b_1^!L)$) deduzida de $u$ (resp. $v$) compondo com o homomorfismo de traço 6.8.7 $a_{2*}a_2^*M\to M$ (resp. $b_{1*}b_1^*L\to L$), ou, o que é equivalente (6.8.6), tomando o produto com $\cl_{Z/Y}(A)$ (resp. $\cl_{Z/X}(B)$). Observemos de passagem que $\cl_{Z/Y}(A)$ (resp. $\cl_{Z/X}(B)$) é a componente em $H^0(A,\Lambda^mN_{A/Z}^{\vee}\otimes\Omega^m_{Z/Y})$ (resp. $H^0(B,\Lambda^nN_{B/Z}^{\vee}\otimes\Omega^n_{Z/X})$) de $\cl_{Z/S}(A)$ (resp. $\cl_{Z/S}(B)$), definida pela projeção $\Omega^m_{Z/S}\to\Omega^m_{Z/Y}$ (resp. $\Omega^n_{Z/S}\to\Omega^n_{Z/X}$) de 6.9.0. Verifica-se facilmente que o produto cup
\[
\langle u^!,v^!\rangle\in\Ext^{m+n}_{\OO_Z}(\OO_C,\Omega^{m+n}_{Z/S})=H^0(C,\Lambda^{m+n}N_{C/Z}^{\vee}\otimes\Omega^{m+n}_{Z/S})
\]
definido em 6.6.3 é dado por
\begin{equation}
\tag{6.9.4}
\langle u^!,v^!\rangle=\langle u,v\rangle\cl_{Z/Y}(A)\cl_{Z/X}(B).
\end{equation}
Além disso, verifica-se que o homomorfismo 6.7.3
\[
\Ext^{m+n}_{\OO_Z}(\OO_C,\Omega^{m+n}_{Z/S})\longrightarrow H^0(S,\OO_S)
\]
é a composta da flecha canônica $\Ext^{m+n}_{\OO_Z}(\OO_C,\Omega^{m+n}_{Z/S})\to H^{m+n}(Z,\Omega^{m+n}_{Z/S})$ e do morfismo de traço $\operatorname{Tr}_{Z/S}:H^{m+n}(Z,\Omega^{m+n}_{Z/S})\to H^0(S,\OO_S)$. A fórmula de Lefschetz--Verdier 6.7.2 fornece, portanto, a relação
\begin{equation}
\tag{6.9.5}
\langle (u^!)_*,(v^!)_*\rangle=\operatorname{Tr}_{Z/S}(\langle u,v\rangle\cl_{Z/Y}(A)\cl_{Z/X}(B)).
\end{equation}
Graças a 6.8.4, o segundo membro pode ser reescrito como um resíduo; em resumo, obtivemos:

